% =============================================================
%  Section 9 -- Discussion
%  Label: sec:discussion
% =============================================================

\section{Discussion}\label{sec:discussion}

\subsection{Transitory Shocks and Rational Forecast Errors}%
\label{sec:transitory}

A natural objection is that the negative forecast errors
reflect rational uncertainty about shock persistence rather
than behavioural overreaction.
If households observe a large meat-price spike but cannot know
ex ante whether it will persist or reverse, a rational agent
who assigns positive probability to persistence will form
expectations that, ex post, overshoot when the shock turns out
to be transitory.
This section shows that the magnitudes rule out this
explanation.

The rational forecast-error coefficient is
$\beta_3^R = \omega(\rho - \hat{\rho})$, where $\omega$ is
the CPI weight on meat, $\rho$ is the true persistence, and
$\hat{\rho}$ is the household's belief about persistence
(Section~\ref{sec:framework}).
Because $\rho$ and $\hat{\rho}$ both lie in $[0,1]$, the
maximum absolute value of the rational coefficient is
$|\beta_3^R| \leq \omega$.
With $\omega \approx 0.05$, the rational forecast-error
coefficient cannot exceed $0.05$ in absolute value regardless
of what the household believes about persistence.

The most generous case for the rational explanation is
$\hat{\rho} = 1$ and $\rho = 0$: the household believes the
shock is permanent but it is fully transitory.
Even then, $\beta_3^R = -0.05$.
The estimated coefficient is $\hat{\beta}_3 = -6.067$ under
region-by-wave fixed effects and $-5.622$ under
province-plus-wave fixed effects---exceeding the rational upper
bound by factors of 121 and 112 respectively.
No intermediate belief $\hat{\rho} \in (0,1)$ can close this
gap.

Table~\ref{tab:calibration_benchmark} makes the comparison
explicit across a grid of persistence assumptions.
At the estimated AR(1) persistence of $\hat{\rho} = 0.480$
(3,689 monthly observations, 31 provinces), a household that
overestimates persistence by the maximum possible amount
produces $\beta_3^R = -0.026$.
The estimated coefficient is 233 times larger in absolute
value.

A subtler version of the objection invokes model uncertainty.
Perhaps households assign a larger CPI weight to meat than the
official weight implies.
To close the gap, households would need to believe that
meat carries a CPI weight of approximately 6 percentage
points---120 times the actual weight.
Even doubling the official weight to $0.10$ leaves the rational
bound 60 times below the estimated coefficient.

Under the diagnostic model of Section~\ref{sec:diagnostic_extension},
the excess forecast error beyond the rational bound identifies
$\theta \cdot \omega$.
The implied $\theta \approx 120$ is large; the ordinal scale
inflates this estimate, so the point value should not be taken
literally.
But the qualitative conclusion is unambiguous: the household
response to meat-price shocks is far beyond what any rational
model with known or uncertain persistence can generate.
The transitory-shock explanation is quantitatively insufficient
under the benchmark, though the exact ratio depends on the
ordinal scaling (Appendix~\ref{app:scaling}).

\subsection{Overreaction Across Demographic Groups}\label{sec:broadbased}

The null demographic heterogeneity results in
Table~\ref{tab:heterogeneity} are informative about the
mechanism driving overreaction.
The effect does not concentrate among low-income or
less-educated households, as a pure expenditure-weight channel
would predict.
Instead, the effect is roughly uniform across the income and
education distribution.

Engel's Law predicts that food's expenditure share declines
with income.
In China, the lowest income quintile allocates 40 to 50 percent
of spending to food, compared with 20 to 25 percent for the
highest quintile.
If overreaction operated through personal budget shares,
low-income households should overreact more by a factor of
roughly two.
The data do not support this prediction: the interaction between
the meat shock and a low-income indicator is small and
insignificant, as are the interactions with education and
urban-rural status.

The absence of demographic gradients is more consistent with
a cognitive salience mechanism.
Under salience, what matters is the visibility of a price
change, not its objective weight in the household's basket.
Pork prices in China are observed frequently at wet markets,
displayed on price boards, and covered extensively in media.
During the ASF episode, pork-price reports were a daily fixture
on television news and social media.
This visibility is common across demographic groups: a
university-educated urban resident and a low-income rural
farmer both encounter pork-price information regularly,
even if their pork expenditure shares differ substantially.

The cross-group similarity in overreaction also has
implications for aggregation.
If overreaction were concentrated among a demographic
minority, aggregate expectation indices would attenuate the
effect by averaging.
Uniform overreaction implies that aggregate measures---such
as the PBoC Urban Depositor Survey---reflect the full
magnitude of the distortion, strengthening the policy
relevance for central banks that rely on survey expectations.

\subsection{External Validity}\label{sec:external}

China's institutional setting amplifies the salience channel
along two dimensions.
First, the official CPI food weight (approximately 28 percent)
is roughly twice the comparable share in high-income economies
(8 percent in the US, 15 percent in the euro area).
Second, pork accounts for approximately 60 percent of total
meat consumption in China---no single protein commodity plays
a comparable role in most OECD countries.
These features suggest that the magnitude documented here may
be an upper bound on what the same design would find in
high-income settings.

But the underlying mechanism---salient relative-price movements
distorting aggregate inflation beliefs---is not specific to China.
\citet{Cavallo2017} document that US households weight
frequently purchased goods more heavily in their inflation
perceptions than official CPI weights imply.
\citet{Coibion2015} find that gasoline-price movements
explain a substantial share of cross-sectional variation
in US household inflation expectations.
\citet{DacuntoHoangPirolliWeber2021} find similar patterns
in European data, and \citet{Malmendier2016} show that
personally experienced inflation episodes have persistent
effects beyond what rational learning models predict.
The cognitive mechanism is portable to other settings where
salient commodity-price shocks generate disproportionate
expectation responses.

\subsection{Policy Implications}\label{sec:policy}

Central banks that use survey-based inflation expectations
as a policy input face a signal-extraction problem of their
own.
If households systematically extrapolate from salient
relative-price movements, raw survey measures overstate the
inflation signal during periods of food-price volatility.
A central bank that tightens in response to elevated survey
expectations driven by a pork-price spike may be responding
to noise rather than signal.

The PBoC publishes quarterly Urban Depositor Survey
expectations as one input to its policy deliberations and
has cited the depositor survey index in quarterly reports
as evidence of expectation anchoring or de-anchoring.
If the index rises mechanically when pork prices spike,
the policy inference is misleading: the rise reflects a
commodity-specific salience effect rather than a generalized
shift in the expected inflation regime.

Adjusting for the commodity-specific salience wedge could
improve the information content of survey expectations.
One approach is to estimate the salience component in real
time using the reduced-form coefficients and observed
meat-price movements, then subtract the predicted contribution
from the raw survey index---analogous to how central banks
use core inflation alongside headline inflation.
The survey design itself could also be improved: decomposing
the question into food and non-food components would allow
direct measurement of the salience wedge.
The Michigan Survey of Consumers has moved in this direction
by asking about expected price changes for specific goods
alongside the general inflation question.

More broadly, central bank communication that explicitly
distinguishes relative-price adjustments from aggregate
inflation trends could dampen the salience channel, though
the effectiveness of such communication for household
expectations (as opposed to financial market expectations)
remains an open question.
